\documentclass[../main]{subfiles}
\begin{document}
\section{Figuras Planas}
\subsection{Aplicaciones del dibujo técnico.}
\newcommand{\ancho}{0.90}
El dibujo técnico es un sistema de comunicación gráfica cuyo propósito es proporcionar información
suficiente para facilitar la interpretación, el análisis, la elaboración de diseños o la resolución de problemas,
por todo esto su finalidad específica es dotar al estudiante de herramientas necesarias para poder
comunicarse gráficamente con objetividad. En la actualidad el Dibujo Técnico se emplea en cualquier
proceso de investigación o proyecto, como lenguaje universal que permite expresar, elaborar e interpretar
información comprensible por cualquier destinatario. Está presente en múltiples situaciones comunicativas
cotidianas, convirtiéndolo en lenguaje básico de comunicación, fiable y objetivo. 
\subsection*{Completar}
\includegraphics[width=0.9\textwidth]{images/04/1.png}

\includegraphics[width=0.9\textwidth]{images/04/2.png}

\includegraphics[width=0.9\textwidth]{images/04/3.png}

\includegraphics[width=0.9\textwidth]{images/04/4.png}

\includegraphics[width=0.9\textwidth]{images/04/5.png}

\includegraphics[width=0.9\textwidth]{images/04/6.png}

\includegraphics[width=0.9\textwidth]{images/04/7.png}

\includegraphics[width=0.9\textwidth]{images/04/8.png}

\newpage

\subsection{Práctica del trazado de líneas} 

La puesta en práctica de las técnicas indicadas es el factor clave para el logro de
resultados. Se sugiere entonces, preparar en una hoja tamaño A4 una serie de puntos
dispuestos como se muestra en las figuras que siguen, separados aproximadamente
cinco milímetros entre si. El objetivo es completar con líneas paralelas a la trazada de
muestra en Figura 3, Figura 4, Figura 5 y Figura 6

\begin{center}
\includegraphics[width=0.6\textwidth]{images/04/9b.png}
\end{center}

El trabajo terminado debiera verse similar a la muestra reducida de la siguiente página.

\includegraphics[width=0.95\textwidth]{images/04/10.png}
\end{document}